%% ch16.tex --- Principal Component Analysis
%% Numbers are computed by verify/ch16_pca.py.
%% C30 is the dataset of the companion lecture's axis-search animation.
%% Figure data (C30 points, digit scatter, reconstructions) come from the same
%% script's output.

\chapter{Principal Component Analysis}
\label{ch:ch16}

\epigraph{The purpose of computing is insight, not numbers.}{---\ Richard Hamming}

\begin{chapterlearning}
After working through this chapter you will be able to:
\begin{enumerate}
  \item Explain why and when to reduce the number of features.
  \item State the principal-components problem as maximising the variance of a
        projection, and derive its solution with a Lagrange multiplier.
  \item Carry out PCA --- centre, scale, decompose, project --- by hand on a
        small dataset, and check the result with the singular value
        decomposition.
  \item Reconstruct data from its principal components, and choose the number of
        components from the variance explained.
  \item Use PCA to visualise high-dimensional data and to compress it.
  \item Explain why PCA is not linear regression, and why it is a poor remedy for
        overfitting.
  \item Say what t-SNE and UMAP do differently.
\end{enumerate}
\end{chapterlearning}

A dataset with fifty features cannot be plotted, and is often more complicated
than it needs to be. Features tend to come in correlated groups --- a car's length
and width, a country's income and energy use, neighbouring pixels in an image ---
so that the data lies close to a surface of much lower dimension than the number
of features. Principal component analysis finds the best flat such surface: a few
new axes along which the data varies most. It is the oldest method in this book,
and still one of the most used.

\section{Reducing the Number of Features}
\label{sec:reduce}

\begin{defbox}{Dimensionality reduction}
\term{Dimensionality reduction} replaces each example $\vx \in \Real{n}$ by a
shorter vector $\mathbf{z} \in \Real{k}$, $k < n$, that retains as much of the
information in $\vx$ as possible. \term{Principal component analysis} (PCA) does
so with a linear map: $\mathbf{z}$ is the coordinates of $\vx$ along $k$ new
orthogonal axes, the \term{principal components}.
\end{defbox}

The uses are three: \emph{visualisation}, by reducing to two or three dimensions
and plotting; \emph{compression}, by storing $k$ numbers instead of $n$; and,
less often now than when computers were slower, \emph{speeding up} a supervised
learner by giving it fewer inputs. PCA is unsupervised: it looks only at the
features, never at a label.

\section{The Idea: Keep the Most Variance}
\label{sec:pca-idea}

Project the data onto a line through its mean. Each point is replaced by its
coordinate along the line, and the information kept is the spread of those
coordinates: a line along which the points are spread out preserves their
differences; a line along which they bunch together destroys them. PCA chooses
the line that maximises the variance of the projections.

\begin{defbox}{Dataset C30}
Thirty points with two strongly correlated features (correlation $0.85$), both
standardised to mean $0$ and variance $1$.
\end{defbox}

\begin{figure}[htbp]
  \centering
  \begin{minipage}[t]{0.48\linewidth}
    \centering
    \begin{tikzpicture}
      \begin{axis}[mlaxis, width=\linewidth, height=\linewidth, axis equal image,
        xmin=-3, xmax=3, ymin=-3, ymax=3, xlabel={$x_1$}, ylabel={$x_2$}]
        \addplot[MLTeal, very thick, domain=-3:3, forget plot] {x};
        \draw[MLGray!60] (axis cs:-1.726,-1.542) -- (axis cs:-1.634,-1.634); \draw[MLGray!60] (axis cs:-1.799,-1.496) -- (axis cs:-1.647,-1.647); \draw[MLGray!60] (axis cs:-1.138,-0.810) -- (axis cs:-0.974,-0.974); \draw[MLGray!60] (axis cs:0.261,-0.470) -- (axis cs:-0.104,-0.104); \draw[MLGray!60] (axis cs:-2.136,-1.390) -- (axis cs:-1.763,-1.763); \draw[MLGray!60] (axis cs:0.438,-0.332) -- (axis cs:0.053,0.053); \draw[MLGray!60] (axis cs:-0.255,-1.012) -- (axis cs:-0.633,-0.633); \draw[MLGray!60] (axis cs:0.653,0.592) -- (axis cs:0.623,0.622); \draw[MLGray!60] (axis cs:0.203,1.050) -- (axis cs:0.627,0.626); \draw[MLGray!60] (axis cs:0.986,1.264) -- (axis cs:1.125,1.125); \draw[MLGray!60] (axis cs:1.187,1.170) -- (axis cs:1.179,1.178); \draw[MLGray!60] (axis cs:-0.188,-0.293) -- (axis cs:-0.240,-0.240); \draw[MLGray!60] (axis cs:0.417,0.835) -- (axis cs:0.626,0.626); \draw[MLGray!60] (axis cs:1.056,2.038) -- (axis cs:1.547,1.547); \draw[MLGray!60] (axis cs:-0.473,-0.701) -- (axis cs:-0.587,-0.587); \draw[MLGray!60] (axis cs:0.661,0.709) -- (axis cs:0.685,0.685); \draw[MLGray!60] (axis cs:0.355,0.534) -- (axis cs:0.445,0.444); \draw[MLGray!60] (axis cs:1.733,1.447) -- (axis cs:1.590,1.590); \draw[MLGray!60] (axis cs:-0.933,-1.014) -- (axis cs:-0.974,-0.974); \draw[MLGray!60] (axis cs:-0.204,-1.093) -- (axis cs:-0.648,-0.648); \draw[MLGray!60] (axis cs:0.087,1.207) -- (axis cs:0.647,0.647); \draw[MLGray!60] (axis cs:0.374,0.360) -- (axis cs:0.367,0.367); \draw[MLGray!60] (axis cs:-1.728,-0.991) -- (axis cs:-1.360,-1.359); \draw[MLGray!60] (axis cs:0.064,0.100) -- (axis cs:0.082,0.082); \draw[MLGray!60] (axis cs:0.683,0.563) -- (axis cs:0.623,0.623); \draw[MLGray!60] (axis cs:-0.088,-0.044) -- (axis cs:-0.066,-0.066); \draw[MLGray!60] (axis cs:0.181,-0.151) -- (axis cs:0.015,0.015); \draw[MLGray!60] (axis cs:1.119,-0.188) -- (axis cs:0.466,0.466); \draw[MLGray!60] (axis cs:-1.275,-1.490) -- (axis cs:-1.383,-1.382); \draw[MLGray!60] (axis cs:1.485,1.146) -- (axis cs:1.316,1.315);
        \addplot[only marks, mark=*, mark size=1.4pt, MLBlue, forget plot] coordinates {(-1.726,-1.542)(-1.799,-1.496)(-1.138,-0.810)(0.261,-0.470)(-2.136,-1.390)(0.438,-0.332)(-0.255,-1.012)(0.653,0.592)(0.203,1.050)(0.986,1.264)(1.187,1.170)(-0.188,-0.293)(0.417,0.835)(1.056,2.038)(-0.473,-0.701)(0.661,0.709)(0.355,0.534)(1.733,1.447)(-0.933,-1.014)(-0.204,-1.093)(0.087,1.207)(0.374,0.360)(-1.728,-0.991)(0.064,0.100)(0.683,0.563)(-0.088,-0.044)(0.181,-0.151)(1.119,-0.188)(-1.275,-1.490)(1.485,1.146)};
      \end{axis}
    \end{tikzpicture}
  \end{minipage}\hfill
  \begin{minipage}[t]{0.48\linewidth}
    \centering
    \begin{tikzpicture}
      \begin{axis}[mlaxis, width=\linewidth, height=\linewidth,
        xmin=0, xmax=180, ymin=0, ymax=2.05, xtick={0,45,90,135,180},
        xlabel={axis angle (degrees)}, ylabel={variance of projections}]
        \addplot[MLBlue, very thick] coordinates {(0,1.0000)(5,1.1470)(10,1.2896)(15,1.4233)(20,1.5442)(25,1.6486)(30,1.7332)(35,1.7956)(40,1.8338)(45,1.8467)(50,1.8338)(55,1.7956)(60,1.7332)(65,1.6486)(70,1.5442)(75,1.4233)(80,1.2896)(85,1.1470)(90,1.0000)(95,0.8530)(100,0.7104)(105,0.5767)(110,0.4558)(115,0.3514)(120,0.2668)(125,0.2044)(130,0.1662)(135,0.1533)(140,0.1662)(145,0.2044)(150,0.2668)(155,0.3514)(160,0.4558)(165,0.5767)(170,0.7104)(175,0.8530)(180,1.0000)};
        \addplot[only marks, mark=*, mark size=2.2pt, MLTeal] coordinates {(45,1.8467)};
        \addplot[only marks, mark=*, mark size=2.2pt, MLAmber] coordinates {(135,0.1533)};
      \end{axis}
    \end{tikzpicture}
  \end{minipage}
  \caption{Left: C30, with the first principal component (teal) and each point's
  projection onto it (grey segments). Right: the variance of the projections as
  the axis turns through $180^\circ$ (the companion lecture animates this search).
  It peaks at $1.8467$ at $45^\circ$ --- the first component --- and is least,
  $0.1533$, at $135^\circ$, the perpendicular axis.}
  \label{fig:c30}
\end{figure}

Two observations from \cref{fig:c30} generalise. The best axis is where the
variance curve peaks, and the worst is perpendicular to it. And the two extreme
variances add to $2$, the total variance of two standardised features: rotating
the axes redistributes variance between them but cannot create or destroy it.

\section{Deriving the Principal Components}
\label{sec:pca-derivation}

\begin{derivation}{The first principal component}
Let the data be centred, $\frac1m\sum_i\xii{i} = \mathbf{0}$, and let
$\Sigma = \frac1m\sum_i\xii{i}\xii{i}{}^{\mathsf T} = \frac1m\Xmat^{\mathsf T}\Xmat$
be its \term{covariance matrix}. The projection of $\xii{i}$ onto a unit vector
$\mathbf{u}$ is $z^{(i)} = \mathbf{u}\cdot\xii{i}$, which has mean $0$ and variance
\[
  \frac1m\sum_i\bigl(\mathbf{u}^{\mathsf T}\xii{i}\bigr)^2
  = \mathbf{u}^{\mathsf T}\Bigl(\frac1m\sum_i\xii{i}\xii{i}{}^{\mathsf T}\Bigr)\mathbf{u}
  = \mathbf{u}^{\mathsf T}\Sigma\,\mathbf{u} .
\]
Maximise it subject to $\mathbf{u}^{\mathsf T}\mathbf{u} = 1$; without the
constraint the variance could be made arbitrarily large by lengthening
$\mathbf{u}$. Introduce a Lagrange multiplier $\lambda$
(\cref{ch:appA}):
\[
  \mathcal{L}(\mathbf{u},\lambda) = \mathbf{u}^{\mathsf T}\Sigma\mathbf{u}
    - \lambda\bigl(\mathbf{u}^{\mathsf T}\mathbf{u} - 1\bigr),
  \qquad
  \nabla_{\mathbf{u}}\mathcal{L} = 2\Sigma\mathbf{u} - 2\lambda\mathbf{u} = \mathbf{0}.
\]
So the optimal $\mathbf{u}$ satisfies
\[
  \boxed{\;\Sigma\,\mathbf{u} = \lambda\,\mathbf{u}\;}
\]
--- it is an \term{eigenvector} of the covariance matrix. Its variance is then
$\mathbf{u}^{\mathsf T}\Sigma\mathbf{u} = \lambda\,\mathbf{u}^{\mathsf T}\mathbf{u}
= \lambda$, the eigenvalue, so the maximum is attained by the eigenvector with the
\emph{largest} eigenvalue.
\end{derivation}

Further components follow the same pattern. The second component maximises the
projected variance subject to being a unit vector \emph{orthogonal} to the first;
repeating the argument with a second multiplier for the orthogonality constraint
shows it is the eigenvector with the second-largest eigenvalue, and so on. Because
$\Sigma$ is symmetric, its eigenvectors can be chosen orthonormal and its
eigenvalues are real and non-negative. The components are therefore
\[
  \mathbf{u}_1, \mathbf{u}_2, \dots, \mathbf{u}_n, \qquad
  \lambda_1 \ge \lambda_2 \ge \dots \ge \lambda_n \ge 0,
\]
with $\lambda_j$ the variance captured by component $j$, and the eigenvalues sum to
the total variance $\operatorname{tr}\Sigma$.

\section{The Algorithm}
\label{sec:pca-algorithm}

\begin{algorithmbox}{Principal component analysis}
\begin{lstlisting}[style=pseudocode]
centre:   subtract each feature's mean (training set)
scale:    divide by each feature's standard deviation,
          if the features have different units
decompose: Sigma = X^T X / m;  eigenvectors u_1..u_n,
           eigenvalues l_1 >= ... >= l_n
          (or SVD X = U S V^T: u_j = row j of V^T,
           l_j = s_j^2 / m)
project:  z_j = x . u_j  for j = 1..k
\end{lstlisting}
The singular value decomposition route avoids forming $\Xmat^{\mathsf T}\Xmat$
explicitly, which is more accurate numerically, and is what libraries use.
\end{algorithmbox}

\begin{workedexample}{Five points by hand}
Take the centred points $(2,1)$, $(1,1)$, $(0,0)$, $(-1,-1)$, $(-2,-1)$. Then
$\sum x_1^2 = 10$, $\sum x_1x_2 = 6$, $\sum x_2^2 = 4$, so
\[
  \Sigma = \frac15\begin{bmatrix}10 & 6\\ 6 & 4\end{bmatrix}
         = \begin{bmatrix}2 & 1.2\\ 1.2 & 0.8\end{bmatrix}.
\]
\textbf{Eigenvalues.} For a $2\times2$ matrix they are
$\bigl(\operatorname{tr}\pm\sqrt{\operatorname{tr}^2 - 4\det}\bigr)/2$; here
$\operatorname{tr} = 2.8$, $\det = 1.6 - 1.44 = 0.16$, and
$\sqrt{7.84 - 0.64} = 2.6833$, so
\[
  \lambda_1 = 2.7416, \qquad \lambda_2 = 0.0584 .
\]
\textbf{Eigenvector.} From $(\Sigma - \lambda_1 I)\mathbf{u} = \mathbf{0}$, the
first row gives $(2 - 2.7416)u_1 + 1.2\,u_2 = 0$, so $u_2/u_1 = 0.6180$ --- the
golden ratio's reciprocal, by a happy accident of these numbers. Normalising,
$\mathbf{u}_1 = (0.8507,\ 0.5257)$.

\textbf{Project.} $z = \vx\cdot\mathbf{u}_1$ gives
$(2.2270,\ 1.3764,\ 0,\ -1.3764,\ -2.2270)$, whose variance is $2.7416$, as the
derivation promised. The first component keeps $\lambda_1/(\lambda_1 +
\lambda_2) = 97.92\%$ of the variance.

\textbf{Check with the SVD.} The singular values of the $5\times2$ data matrix are
$3.7025$ and $0.5402$; their squares divided by $m = 5$ are $2.7416$ and $0.0584$,
and the first right singular vector is $(-0.8507,\ -0.5257)$ --- the same
direction, with the opposite sign. The sign of an eigenvector is arbitrary;
implementations may return either.
\end{workedexample}

\begin{pitfall}{PCA without centring or scaling}
The covariance formula assumes centred data. Without centring, the first
``component'' points from the origin towards the data's mean rather than along its
spread. Without scaling, features measured in large units dominate: a length in
millimetres has a variance a million times that of the same length in metres, and
PCA will declare it the most important direction for no reason but the choice of
unit. Centre always; scale whenever the features are in different units. Compute
the means and standard deviations on the training set, and reuse them for new
data.
\end{pitfall}

\begin{researchspotlight}{Pearson (1901) --- On lines and planes of closest fit}
\textbf{The problem.} Given a cloud of points in many dimensions, find the line ---
or plane --- that fits it best, where every coordinate is measured with error and
none is singled out as the one to be predicted.

\textbf{The idea.} Minimise the sum of squared \emph{perpendicular} distances from
the points to the line, rather than the vertical distances of regression. Pearson
showed that the best line passes through the centroid along the direction of
greatest spread, and that the best plane contains it.

\textbf{The result.} The first statement of principal component analysis, with
geometric arguments and worked figures. Hotelling (1933) later gave the
algebraic, eigenvector formulation and the name.

\textbf{Why it is here.} Pearson's ``closest fit'' and the ``maximum variance'' of
\cref{sec:pca-derivation} are the same line --- the reconstruction view of
\cref{sec:reconstruction} --- and his contrast with regression is the pitfall of
\cref{sec:pca-misuse}.
\parencite{pearson1901}
\end{researchspotlight}

\section{Reconstruction and Choosing \texorpdfstring{$k$}{k}}
\label{sec:reconstruction}

From the $k$ coordinates $z_j = \vx\cdot\mathbf{u}_j$, the best approximation of
$\vx$ is
\[
  \hat\vx = \sum_{j=1}^{k} z_j\,\mathbf{u}_j .
\]
For the first of the five points, $z = 2.2270$ gives $\hat\vx = 2.2270\,(0.8507,\
0.5257) = (1.8944,\ 1.1708)$, against the original $(2,\ 1)$. Averaged over the
points, the squared reconstruction error is $0.0584 = \lambda_2$: the variance
PCA discards is exactly the error it makes. In general,
\[
  \frac1m\sum_i\bigl\|\xii{i} - \hat\vx^{(i)}\bigr\|^2 = \sum_{j=k+1}^{n}\lambda_j ,
\]
so maximising the variance kept and minimising the reconstruction error are the
same problem --- Pearson's view and Hotelling's, reconciled.

\begin{defbox}{Variance explained}
The fraction of the total variance kept by the first $k$ components is
\[
  \frac{\lambda_1 + \dots + \lambda_k}{\lambda_1 + \dots + \lambda_n} .
\]
A common rule is to choose the smallest $k$ that keeps $90\%$, $95\%$ or $99\%$,
depending on how much loss the application tolerates.
\end{defbox}

\begin{workedexample}{The digits: 64 features}
Each of the $1797$ handwritten digits in \tool{scikit-learn}'s \code{load\_digits}
is an $8\times8$ greyscale image, a point in $\Real{64}$. With brightness scaled to
$[0,1]$:

\begin{center}
\small
\renewcommand{\arraystretch}{1.12}
\begin{tabular}{@{}rrrr@{}}
\toprule
$k$ & variance kept & mean squared error per pixel & storage\\
\midrule
$2$  & $28.5\%$ & $0.0524$ & $3.1\%$\\
$5$  & $54.5\%$ & $0.0334$ & $7.8\%$\\
$10$ & $73.8\%$ & $0.0192$ & $15.6\%$\\
$15$ & $83.5\%$ & $0.0121$ & $23.4\%$\\
$30$ & $95.9\%$ & $0.0030$ & $46.9\%$\\
\bottomrule
\end{tabular}
\end{center}

Keeping $90\%$ of the variance needs $21$ components, $95\%$ needs $29$, and $99\%$
needs $41$ --- the data is genuinely spread over many directions, but a third of
the original size keeps nearly everything.
\end{workedexample}

\begin{figure}[htbp]
  \centering
  \begin{tikzpicture}[x=2.6mm,y=2.6mm]
    \fill[black!0] (0,0) rectangle ++(1,1);
    \fill[black!0] (1,0) rectangle ++(1,1);
    \fill[black!44] (2,0) rectangle ++(1,1);
    \fill[black!94] (3,0) rectangle ++(1,1);
    \fill[black!81] (4,0) rectangle ++(1,1);
    \fill[black!6] (5,0) rectangle ++(1,1);
    \fill[black!0] (6,0) rectangle ++(1,1);
    \fill[black!0] (7,0) rectangle ++(1,1);
    \fill[black!0] (0,-1) rectangle ++(1,1);
    \fill[black!50] (1,-1) rectangle ++(1,1);
    \fill[black!81] (2,-1) rectangle ++(1,1);
    \fill[black!38] (3,-1) rectangle ++(1,1);
    \fill[black!94] (4,-1) rectangle ++(1,1);
    \fill[black!25] (5,-1) rectangle ++(1,1);
    \fill[black!0] (6,-1) rectangle ++(1,1);
    \fill[black!0] (7,-1) rectangle ++(1,1);
    \fill[black!0] (0,-2) rectangle ++(1,1);
    \fill[black!12] (1,-2) rectangle ++(1,1);
    \fill[black!6] (2,-2) rectangle ++(1,1);
    \fill[black!81] (3,-2) rectangle ++(1,1);
    \fill[black!81] (4,-2) rectangle ++(1,1);
    \fill[black!0] (5,-2) rectangle ++(1,1);
    \fill[black!0] (6,-2) rectangle ++(1,1);
    \fill[black!0] (7,-2) rectangle ++(1,1);
    \fill[black!0] (0,-3) rectangle ++(1,1);
    \fill[black!0] (1,-3) rectangle ++(1,1);
    \fill[black!12] (2,-3) rectangle ++(1,1);
    \fill[black!94] (3,-3) rectangle ++(1,1);
    \fill[black!69] (4,-3) rectangle ++(1,1);
    \fill[black!6] (5,-3) rectangle ++(1,1);
    \fill[black!0] (6,-3) rectangle ++(1,1);
    \fill[black!0] (7,-3) rectangle ++(1,1);
    \fill[black!0] (0,-4) rectangle ++(1,1);
    \fill[black!0] (1,-4) rectangle ++(1,1);
    \fill[black!0] (2,-4) rectangle ++(1,1);
    \fill[black!6] (3,-4) rectangle ++(1,1);
    \fill[black!75] (4,-4) rectangle ++(1,1);
    \fill[black!75] (5,-4) rectangle ++(1,1);
    \fill[black!6] (6,-4) rectangle ++(1,1);
    \fill[black!0] (7,-4) rectangle ++(1,1);
    \fill[black!0] (0,-5) rectangle ++(1,1);
    \fill[black!0] (1,-5) rectangle ++(1,1);
    \fill[black!0] (2,-5) rectangle ++(1,1);
    \fill[black!0] (3,-5) rectangle ++(1,1);
    \fill[black!6] (4,-5) rectangle ++(1,1);
    \fill[black!62] (5,-5) rectangle ++(1,1);
    \fill[black!50] (6,-5) rectangle ++(1,1);
    \fill[black!0] (7,-5) rectangle ++(1,1);
    \fill[black!0] (0,-6) rectangle ++(1,1);
    \fill[black!0] (1,-6) rectangle ++(1,1);
    \fill[black!50] (2,-6) rectangle ++(1,1);
    \fill[black!25] (3,-6) rectangle ++(1,1);
    \fill[black!31] (4,-6) rectangle ++(1,1);
    \fill[black!88] (5,-6) rectangle ++(1,1);
    \fill[black!56] (6,-6) rectangle ++(1,1);
    \fill[black!0] (7,-6) rectangle ++(1,1);
    \fill[black!0] (0,-7) rectangle ++(1,1);
    \fill[black!0] (1,-7) rectangle ++(1,1);
    \fill[black!44] (2,-7) rectangle ++(1,1);
    \fill[black!81] (3,-7) rectangle ++(1,1);
    \fill[black!81] (4,-7) rectangle ++(1,1);
    \fill[black!56] (5,-7) rectangle ++(1,1);
    \fill[black!0] (6,-7) rectangle ++(1,1);
    \fill[black!0] (7,-7) rectangle ++(1,1);
    \draw[MLGray!60] (0,-7) rectangle ++(8,8);
    \node[font=\small] at (4,-8.3) {original};
    \fill[black!0] (11,0) rectangle ++(1,1);
    \fill[black!6] (12,0) rectangle ++(1,1);
    \fill[black!56] (13,0) rectangle ++(1,1);
    \fill[black!88] (14,0) rectangle ++(1,1);
    \fill[black!75] (15,0) rectangle ++(1,1);
    \fill[black!44] (16,0) rectangle ++(1,1);
    \fill[black!6] (17,0) rectangle ++(1,1);
    \fill[black!0] (18,0) rectangle ++(1,1);
    \fill[black!0] (11,-1) rectangle ++(1,1);
    \fill[black!25] (12,-1) rectangle ++(1,1);
    \fill[black!88] (13,-1) rectangle ++(1,1);
    \fill[black!62] (14,-1) rectangle ++(1,1);
    \fill[black!69] (15,-1) rectangle ++(1,1);
    \fill[black!69] (16,-1) rectangle ++(1,1);
    \fill[black!12] (17,-1) rectangle ++(1,1);
    \fill[black!0] (18,-1) rectangle ++(1,1);
    \fill[black!0] (11,-2) rectangle ++(1,1);
    \fill[black!25] (12,-2) rectangle ++(1,1);
    \fill[black!56] (13,-2) rectangle ++(1,1);
    \fill[black!19] (14,-2) rectangle ++(1,1);
    \fill[black!56] (15,-2) rectangle ++(1,1);
    \fill[black!62] (16,-2) rectangle ++(1,1);
    \fill[black!6] (17,-2) rectangle ++(1,1);
    \fill[black!0] (18,-2) rectangle ++(1,1);
    \fill[black!0] (11,-3) rectangle ++(1,1);
    \fill[black!12] (12,-3) rectangle ++(1,1);
    \fill[black!38] (13,-3) rectangle ++(1,1);
    \fill[black!56] (14,-3) rectangle ++(1,1);
    \fill[black!75] (15,-3) rectangle ++(1,1);
    \fill[black!50] (16,-3) rectangle ++(1,1);
    \fill[black!12] (17,-3) rectangle ++(1,1);
    \fill[black!0] (18,-3) rectangle ++(1,1);
    \fill[black!0] (11,-4) rectangle ++(1,1);
    \fill[black!0] (12,-4) rectangle ++(1,1);
    \fill[black!12] (13,-4) rectangle ++(1,1);
    \fill[black!38] (14,-4) rectangle ++(1,1);
    \fill[black!50] (15,-4) rectangle ++(1,1);
    \fill[black!56] (16,-4) rectangle ++(1,1);
    \fill[black!19] (17,-4) rectangle ++(1,1);
    \fill[black!0] (18,-4) rectangle ++(1,1);
    \fill[black!0] (11,-5) rectangle ++(1,1);
    \fill[black!0] (12,-5) rectangle ++(1,1);
    \fill[black!12] (13,-5) rectangle ++(1,1);
    \fill[black!12] (14,-5) rectangle ++(1,1);
    \fill[black!19] (15,-5) rectangle ++(1,1);
    \fill[black!56] (16,-5) rectangle ++(1,1);
    \fill[black!31] (17,-5) rectangle ++(1,1);
    \fill[black!0] (18,-5) rectangle ++(1,1);
    \fill[black!0] (11,-6) rectangle ++(1,1);
    \fill[black!6] (12,-6) rectangle ++(1,1);
    \fill[black!56] (13,-6) rectangle ++(1,1);
    \fill[black!50] (14,-6) rectangle ++(1,1);
    \fill[black!44] (15,-6) rectangle ++(1,1);
    \fill[black!75] (16,-6) rectangle ++(1,1);
    \fill[black!31] (17,-6) rectangle ++(1,1);
    \fill[black!0] (18,-6) rectangle ++(1,1);
    \fill[black!0] (11,-7) rectangle ++(1,1);
    \fill[black!0] (12,-7) rectangle ++(1,1);
    \fill[black!56] (13,-7) rectangle ++(1,1);
    \fill[black!88] (14,-7) rectangle ++(1,1);
    \fill[black!81] (15,-7) rectangle ++(1,1);
    \fill[black!56] (16,-7) rectangle ++(1,1);
    \fill[black!19] (17,-7) rectangle ++(1,1);
    \fill[black!6] (18,-7) rectangle ++(1,1);
    \draw[MLGray!60] (11,-7) rectangle ++(8,8);
    \node[font=\small] at (15,-8.3) {$k = 2$};
    \fill[black!0] (22,0) rectangle ++(1,1);
    \fill[black!6] (23,0) rectangle ++(1,1);
    \fill[black!50] (24,0) rectangle ++(1,1);
    \fill[black!75] (25,0) rectangle ++(1,1);
    \fill[black!62] (26,0) rectangle ++(1,1);
    \fill[black!25] (27,0) rectangle ++(1,1);
    \fill[black!0] (28,0) rectangle ++(1,1);
    \fill[black!0] (29,0) rectangle ++(1,1);
    \fill[black!0] (22,-1) rectangle ++(1,1);
    \fill[black!25] (23,-1) rectangle ++(1,1);
    \fill[black!75] (24,-1) rectangle ++(1,1);
    \fill[black!62] (25,-1) rectangle ++(1,1);
    \fill[black!75] (26,-1) rectangle ++(1,1);
    \fill[black!44] (27,-1) rectangle ++(1,1);
    \fill[black!0] (28,-1) rectangle ++(1,1);
    \fill[black!0] (29,-1) rectangle ++(1,1);
    \fill[black!0] (22,-2) rectangle ++(1,1);
    \fill[black!25] (23,-2) rectangle ++(1,1);
    \fill[black!50] (24,-2) rectangle ++(1,1);
    \fill[black!56] (25,-2) rectangle ++(1,1);
    \fill[black!81] (26,-2) rectangle ++(1,1);
    \fill[black!38] (27,-2) rectangle ++(1,1);
    \fill[black!0] (28,-2) rectangle ++(1,1);
    \fill[black!0] (29,-2) rectangle ++(1,1);
    \fill[black!0] (22,-3) rectangle ++(1,1);
    \fill[black!6] (23,-3) rectangle ++(1,1);
    \fill[black!38] (24,-3) rectangle ++(1,1);
    \fill[black!88] (25,-3) rectangle ++(1,1);
    \fill[black!94] (26,-3) rectangle ++(1,1);
    \fill[black!31] (27,-3) rectangle ++(1,1);
    \fill[black!0] (28,-3) rectangle ++(1,1);
    \fill[black!0] (29,-3) rectangle ++(1,1);
    \fill[black!0] (22,-4) rectangle ++(1,1);
    \fill[black!0] (23,-4) rectangle ++(1,1);
    \fill[black!0] (24,-4) rectangle ++(1,1);
    \fill[black!44] (25,-4) rectangle ++(1,1);
    \fill[black!62] (26,-4) rectangle ++(1,1);
    \fill[black!50] (27,-4) rectangle ++(1,1);
    \fill[black!6] (28,-4) rectangle ++(1,1);
    \fill[black!0] (29,-4) rectangle ++(1,1);
    \fill[black!0] (22,-5) rectangle ++(1,1);
    \fill[black!0] (23,-5) rectangle ++(1,1);
    \fill[black!0] (24,-5) rectangle ++(1,1);
    \fill[black!12] (25,-5) rectangle ++(1,1);
    \fill[black!25] (26,-5) rectangle ++(1,1);
    \fill[black!69] (27,-5) rectangle ++(1,1);
    \fill[black!38] (28,-5) rectangle ++(1,1);
    \fill[black!0] (29,-5) rectangle ++(1,1);
    \fill[black!0] (22,-6) rectangle ++(1,1);
    \fill[black!6] (23,-6) rectangle ++(1,1);
    \fill[black!50] (24,-6) rectangle ++(1,1);
    \fill[black!38] (25,-6) rectangle ++(1,1);
    \fill[black!50] (26,-6) rectangle ++(1,1);
    \fill[black!100] (27,-6) rectangle ++(1,1);
    \fill[black!56] (28,-6) rectangle ++(1,1);
    \fill[black!6] (29,-6) rectangle ++(1,1);
    \fill[black!0] (22,-7) rectangle ++(1,1);
    \fill[black!0] (23,-7) rectangle ++(1,1);
    \fill[black!50] (24,-7) rectangle ++(1,1);
    \fill[black!81] (25,-7) rectangle ++(1,1);
    \fill[black!100] (26,-7) rectangle ++(1,1);
    \fill[black!94] (27,-7) rectangle ++(1,1);
    \fill[black!31] (28,-7) rectangle ++(1,1);
    \fill[black!6] (29,-7) rectangle ++(1,1);
    \draw[MLGray!60] (22,-7) rectangle ++(8,8);
    \node[font=\small] at (26,-8.3) {$k = 5$};
    \fill[black!0] (33,0) rectangle ++(1,1);
    \fill[black!6] (34,0) rectangle ++(1,1);
    \fill[black!50] (35,0) rectangle ++(1,1);
    \fill[black!81] (36,0) rectangle ++(1,1);
    \fill[black!75] (37,0) rectangle ++(1,1);
    \fill[black!25] (38,0) rectangle ++(1,1);
    \fill[black!0] (39,0) rectangle ++(1,1);
    \fill[black!0] (40,0) rectangle ++(1,1);
    \fill[black!0] (33,-1) rectangle ++(1,1);
    \fill[black!25] (34,-1) rectangle ++(1,1);
    \fill[black!62] (35,-1) rectangle ++(1,1);
    \fill[black!56] (36,-1) rectangle ++(1,1);
    \fill[black!81] (37,-1) rectangle ++(1,1);
    \fill[black!19] (38,-1) rectangle ++(1,1);
    \fill[black!0] (39,-1) rectangle ++(1,1);
    \fill[black!0] (40,-1) rectangle ++(1,1);
    \fill[black!0] (33,-2) rectangle ++(1,1);
    \fill[black!6] (34,-2) rectangle ++(1,1);
    \fill[black!12] (35,-2) rectangle ++(1,1);
    \fill[black!56] (36,-2) rectangle ++(1,1);
    \fill[black!88] (37,-2) rectangle ++(1,1);
    \fill[black!0] (38,-2) rectangle ++(1,1);
    \fill[black!0] (39,-2) rectangle ++(1,1);
    \fill[black!0] (40,-2) rectangle ++(1,1);
    \fill[black!0] (33,-3) rectangle ++(1,1);
    \fill[black!0] (34,-3) rectangle ++(1,1);
    \fill[black!12] (35,-3) rectangle ++(1,1);
    \fill[black!81] (36,-3) rectangle ++(1,1);
    \fill[black!88] (37,-3) rectangle ++(1,1);
    \fill[black!6] (38,-3) rectangle ++(1,1);
    \fill[black!0] (39,-3) rectangle ++(1,1);
    \fill[black!0] (40,-3) rectangle ++(1,1);
    \fill[black!0] (33,-4) rectangle ++(1,1);
    \fill[black!0] (34,-4) rectangle ++(1,1);
    \fill[black!0] (35,-4) rectangle ++(1,1);
    \fill[black!25] (36,-4) rectangle ++(1,1);
    \fill[black!50] (37,-4) rectangle ++(1,1);
    \fill[black!69] (38,-4) rectangle ++(1,1);
    \fill[black!19] (39,-4) rectangle ++(1,1);
    \fill[black!0] (40,-4) rectangle ++(1,1);
    \fill[black!0] (33,-5) rectangle ++(1,1);
    \fill[black!0] (34,-5) rectangle ++(1,1);
    \fill[black!0] (35,-5) rectangle ++(1,1);
    \fill[black!0] (36,-5) rectangle ++(1,1);
    \fill[black!0] (37,-5) rectangle ++(1,1);
    \fill[black!69] (38,-5) rectangle ++(1,1);
    \fill[black!50] (39,-5) rectangle ++(1,1);
    \fill[black!0] (40,-5) rectangle ++(1,1);
    \fill[black!0] (33,-6) rectangle ++(1,1);
    \fill[black!6] (34,-6) rectangle ++(1,1);
    \fill[black!38] (35,-6) rectangle ++(1,1);
    \fill[black!31] (36,-6) rectangle ++(1,1);
    \fill[black!44] (37,-6) rectangle ++(1,1);
    \fill[black!88] (38,-6) rectangle ++(1,1);
    \fill[black!44] (39,-6) rectangle ++(1,1);
    \fill[black!0] (40,-6) rectangle ++(1,1);
    \fill[black!0] (33,-7) rectangle ++(1,1);
    \fill[black!6] (34,-7) rectangle ++(1,1);
    \fill[black!56] (35,-7) rectangle ++(1,1);
    \fill[black!88] (36,-7) rectangle ++(1,1);
    \fill[black!100] (37,-7) rectangle ++(1,1);
    \fill[black!56] (38,-7) rectangle ++(1,1);
    \fill[black!12] (39,-7) rectangle ++(1,1);
    \fill[black!0] (40,-7) rectangle ++(1,1);
    \draw[MLGray!60] (33,-7) rectangle ++(8,8);
    \node[font=\small] at (37,-8.3) {$k = 15$};
  \end{tikzpicture}
  \caption{One handwritten $3$ (left) and its reconstructions from $2$, $5$ and
  $15$ principal components, drawn as $8\times8$ grids of grey cells. Even two
  components give a recognisable outline, because the reconstruction starts from
  the average digit; the strokes sharpen as components are added. Squared error per
  pixel: $0.0423$, $0.0235$ and $0.0097$.}
  \label{fig:recon}
\end{figure}

\begin{figure}[htbp]
  \centering
  \begin{tikzpicture}
    \begin{axis}[mlaxis, width=0.66\linewidth, height=4.6cm,
      xlabel={number of components $k$}, ylabel={variance kept},
      xmin=0, xmax=41, ymin=0, ymax=1.02]
      \addplot[MLAmber, very thick, mark=*, mark size=0.9pt] coordinates {(1,0.1489)(2,0.2851)(3,0.4030)(4,0.4871)(5,0.5450)(6,0.5941)(7,0.6373)(8,0.6739)(9,0.7074)(10,0.7382)(11,0.7620)(12,0.7847)(13,0.8029)(14,0.8206)(15,0.8353)(16,0.8494)(17,0.8626)(18,0.8751)(19,0.8852)(20,0.8943)(21,0.9032)(22,0.9112)(23,0.9188)(24,0.9261)(25,0.9330)(26,0.9390)(27,0.9447)(28,0.9499)(29,0.9548)(30,0.9591)(31,0.9628)(32,0.9664)(33,0.9697)(34,0.9730)(35,0.9761)(36,0.9790)(37,0.9816)(38,0.9839)(39,0.9861)(40,0.9882)};
      \addplot[MLGray, dashed, domain=0:41, samples=2] {0.9};
      \node[font=\tiny, MLGray, anchor=south east] at (axis cs:40,0.9) {$90\%$};
    \end{axis}
  \end{tikzpicture}
  \caption{Cumulative variance explained for the $64$-dimensional digits. The curve
  rises steeply and then slowly; $21$ components reach $90\%$.}
  \label{fig:cumvar}
\end{figure}

\section{PCA for Visualisation}
\label{sec:pca-vis}

Reduce to two components and plot. Take the digits $0$, $1$ and $7$ only ---
$539$ images --- and fit PCA without using the labels.

\begin{figure}[htbp]
  \centering
  \begin{tikzpicture}
    \begin{axis}[mlaxis, width=0.72\linewidth, height=5.6cm,
      xlabel={first component $z_1$}, ylabel={second component $z_2$},
      legend style={font=\scriptsize, at={(0.02,0.98)}, anchor=north west}]
      \addplot[only marks, mark=*, mark size=1.4pt, MLBlue] coordinates {(-1.557,0.171)(-1.678,-0.264)(-1.190,0.016)(-1.758,-0.231)(-1.543,0.088)(-1.680,0.017)(-1.479,-0.159)(-1.796,0.015)(-1.258,-0.350)(-1.791,0.368)(-1.560,0.280)(-1.790,-0.218)(-2.056,-0.348)(-1.528,-0.291)(-0.810,0.294)(-1.954,-0.275)(-1.846,0.073)(-1.681,0.107)(-1.551,0.578)(-1.770,0.443)(-1.697,-0.189)(-1.573,-0.224)(-1.642,-0.241)(-1.668,-0.650)(-1.815,-0.414)(-0.838,0.786)(-1.775,-0.098)(-1.648,0.421)(-1.826,-0.336)(-2.034,-0.186)(-2.050,-0.017)(-1.679,0.256)(-2.092,0.093)(-1.403,0.117)(-1.726,-0.367)(-1.781,-0.243)(-1.909,0.091)(-1.848,-0.531)(-1.786,-0.184)(-1.738,0.228)};
      \addlegendentry{digit 0}
      \addplot[only marks, mark=square*, mark size=1.4pt, MLRose] coordinates {(1.387,0.658)(0.491,1.758)(1.316,1.228)(1.458,1.073)(0.657,1.783)(1.560,1.024)(0.294,1.646)(1.206,0.663)(0.445,0.944)(0.920,1.313)(0.240,1.695)(1.042,0.783)(1.178,0.893)(1.247,0.862)(1.273,1.026)(1.365,0.856)(0.530,1.099)(0.556,1.245)(0.889,1.004)(1.504,0.703)(0.692,0.699)(0.530,1.425)(0.210,1.751)(0.653,0.751)(0.285,1.848)(1.501,1.198)(0.012,0.602)(1.105,0.968)(-0.506,0.438)(-0.766,0.164)(0.423,1.313)(-0.079,0.306)(0.045,0.952)(1.481,0.627)(0.012,0.306)(0.848,0.744)(1.064,0.995)(1.454,1.131)(0.018,0.827)(0.991,0.835)};
      \addlegendentry{digit 1}
      \addplot[only marks, mark=triangle*, mark size=1.7pt, MLTeal] coordinates {(0.749,-0.473)(0.854,-1.002)(0.651,-1.455)(0.888,-0.458)(0.580,-1.371)(0.574,-1.000)(0.753,-0.922)(0.928,-1.310)(1.060,-1.135)(1.246,-0.687)(1.036,-1.271)(0.882,-0.689)(0.825,-1.315)(0.584,-0.940)(1.198,-0.195)(0.965,-0.292)(0.876,-1.103)(0.256,-0.919)(0.877,-1.525)(0.705,-1.314)(0.351,-0.723)(1.068,-0.686)(0.891,-0.805)(0.829,-0.742)(1.188,-1.263)(0.212,-1.009)(1.002,-1.012)(0.968,-1.011)(0.784,-1.341)(0.387,-0.585)(1.337,-0.092)(1.127,-1.058)(0.595,-0.762)(1.116,-0.502)(1.344,-0.564)(0.998,-0.535)(0.531,-1.106)(0.919,-0.863)(0.438,-1.136)(0.703,-0.539)};
      \addlegendentry{digit 7}
    \end{axis}
  \end{tikzpicture}
  \caption{$40$ images of each digit, projected onto the first two principal
  components of the $539$ images of $0$, $1$ and $7$. The components were found
  without labels, and keep only $47.3\%$ of the variance, yet the three digits
  form three clear groups --- the colours are added afterwards.}
  \label{fig:digits-pca}
\end{figure}

A two-dimensional PCA plot is a standard first look at any new dataset: it shows
clusters, outliers and gradients that no table of numbers reveals. Its axes are
combinations of all the original features, and part of the analysis is to inspect
which features weigh most in each.

\section{What PCA Is Not}
\label{sec:pca-misuse}

\subsection{It is not linear regression}

Both fit a line through a cloud of points, and they fit different lines.
Regression predicts one designated variable from the others and minimises
\emph{vertical} distances. PCA treats all features symmetrically and minimises
\emph{perpendicular} distances. On the five points of \cref{sec:pca-algorithm}, the
least-squares regression of $x_2$ on $x_1$ has slope $0.6000$; the first principal
component has slope $0.6180$. The difference grows with the noise.

\begin{figure}[htbp]
  \centering
  \begin{tikzpicture}
    \begin{axis}[mlaxis, width=0.62\linewidth, height=5.2cm, axis equal image,
      xmin=-2.6, xmax=2.6, ymin=-1.8, ymax=1.8, xlabel={$x_1$}, ylabel={$x_2$},
      legend style={font=\scriptsize, at={(0.02,0.98)}, anchor=north west}]
      \addplot[MLAmber, thick, dashed, domain=-2.6:2.6] {0.6*x};
      \addlegendentry{regression, slope $0.6000$}
      \addplot[MLTeal, very thick, domain=-2.6:2.6] {0.6180*x};
      \addlegendentry{PCA, slope $0.6180$}
      \draw[MLAmber!70] (axis cs:2,1) -- (axis cs:2,1.2);
      \draw[MLAmber!70] (axis cs:-2,-1) -- (axis cs:-2,-1.2);
      \draw[MLAmber!70] (axis cs:1,1) -- (axis cs:1,0.6);
      \draw[MLAmber!70] (axis cs:-1,-1) -- (axis cs:-1,-0.6);
      \draw[MLTeal!70] (axis cs:2,1) -- (axis cs:1.8944,1.1708);
      \draw[MLTeal!70] (axis cs:1,1) -- (axis cs:1.1708,0.7236);
      \draw[MLTeal!70] (axis cs:-1,-1) -- (axis cs:-1.1708,-0.7236);
      \draw[MLTeal!70] (axis cs:-2,-1) -- (axis cs:-1.8944,-1.1708);
      \addplot[only marks, mark=*, mark size=2pt, MLBlue, forget plot]
        coordinates {(2,1)(1,1)(0,0)(-1,-1)(-2,-1)};
    \end{axis}
  \end{tikzpicture}
  \caption{Regression and PCA on the same five points. Regression (dashed) minimises
  the vertical distances (amber segments); PCA (solid) minimises the perpendicular
  ones (teal segments). Only PCA treats $x_1$ and $x_2$ alike.}
  \label{fig:pca-vs-reg}
\end{figure}

\subsection{It is not a cure for overfitting}

It is tempting to reduce the features with PCA before training, to give the model
fewer parameters and so less room to overfit. But PCA chooses the directions of
greatest variance without looking at the labels, and the direction that best
separates the classes may have little variance.

\begin{workedexample}{PCA versus regularization on the digits}
Train ten-class logistic regression on $200$ digit images --- few enough that
overfitting threatens --- and test on $899$:

\begin{center}
\small
\renewcommand{\arraystretch}{1.12}
\begin{tabular}{@{}lrr@{}}
\toprule
model & train accuracy & test accuracy\\
\midrule
all $64$ pixels, almost no penalty ($C = 10^4$) & $1.0000$ & $0.9121$\\
all $64$ pixels, penalty $C = 1$                 & $1.0000$ & $\mathbf{0.9155}$\\
all $64$ pixels, strong penalty $C = 0.1$        & $0.9300$ & $0.8843$\\
PCA to $5$ components, no penalty                & $0.8700$ & $0.8120$\\
PCA to $10$ components, no penalty               & $1.0000$ & $0.8676$\\
PCA to $20$ components, no penalty               & $1.0000$ & $0.9121$\\
\bottomrule
\end{tabular}
\end{center}

(\tool{scikit-learn}'s $C$ is the inverse of the regularization strength,
$C = 1/\lambda$.) Ten components throw away enough information to cost almost
five percentage points; twenty recover the unregularized model's accuracy but no
more. The best model keeps every pixel and controls overfitting with the L2 penalty
of \cref{ch:ch07}.
\end{workedexample}

\begin{pitfall}{Using PCA to prevent overfitting}
Use regularization, more data, or a simpler model class for overfitting. Use PCA
for visualisation and compression --- and if it is used to speed up training,
choose $k$ by validation performance of the downstream model, not by variance
explained alone.
\end{pitfall}

\section{Beyond PCA: t-SNE and UMAP}
\label{sec:tsne}

PCA is linear: it can only rotate and project. Data lying on a curved surface ---
a spiral, or the many distorted versions of a handwritten digit --- is not well
summarised by any flat projection. Two non-linear methods are widely used for
visualisation instead.

\term{t-SNE} places points in two dimensions so that each point's nearest
neighbours in the original space remain near it, and \term{UMAP} pursues a similar
goal with a different mathematical foundation and greater speed. Both reveal
clusters that PCA misses. Both also distort: distances between clusters and the
sizes of clusters in their plots carry little meaning, and the result changes with
their hyperparameters and random seed. Use them to look, not to measure; use PCA
when the axes must mean something or when new data must be projected the same way.

\begin{researchspotlight}{van der Maaten \& Hinton (2008) --- Visualizing data using t-SNE}
\textbf{The problem.} Existing methods for visualising high-dimensional data in two
dimensions either preserved global structure and lost local neighbourhoods (like
PCA), or crowded points together in the centre of the map.

\textbf{The idea.} Convert distances in the original space into probabilities that
points are neighbours, do the same in two dimensions using a heavy-tailed
Student-$t$ distribution, and move the two-dimensional points to make the two sets
of probabilities match. The heavy tail relieves the crowding.

\textbf{The result.} Maps of handwritten digits, faces and text in which the
classes separate strikingly; t-SNE became the default visualisation for
high-dimensional data, including the internal representations of neural networks.

\textbf{Why it is here.} It is the non-linear alternative of \cref{sec:tsne}, and a
lesson in reading plots: its clusters are real, but their spacing is not.
\parencite{vandermaaten2008}
\end{researchspotlight}

\begin{labbox}{Lab 16 --- PCA and data visualisation}
Implement PCA from scratch: centre, form $\Sigma$, compute its eigenvectors with
\code{np.linalg.eigh}, sort by eigenvalue, and project. Reproduce the five-point
example of \cref{sec:pca-algorithm} and confirm it against \code{np.linalg.svd}
and \code{sklearn.decomposition.PCA}.

Then apply PCA to the digits: plot the cumulative variance explained, reconstruct
one image from $2$, $5$ and $15$ components, and draw the two-dimensional map of
\cref{fig:digits-pca}. Finally fit t-SNE to the same images and compare the two
maps.
\end{labbox}

\section*{Chapter Summary}
\addcontentsline{toc}{section}{Chapter Summary}

\begin{itemize}
  \item PCA replaces $n$ correlated features by $k$ coordinates along the
        directions of greatest variance, for visualisation and compression.
  \item The variance of a projection onto unit $\mathbf{u}$ is
        $\mathbf{u}^{\mathsf T}\Sigma\mathbf{u}$; maximising it under
        $\|\mathbf{u}\| = 1$ gives $\Sigma\mathbf{u} = \lambda\mathbf{u}$. The
        components are the eigenvectors of the covariance matrix in order of
        eigenvalue, and each eigenvalue is the variance its component keeps.
  \item Algorithm: centre, scale if units differ, decompose (eigen or SVD),
        project. On five points, $\lambda_1 = 2.7416$ keeps $97.92\%$ of the
        variance along $\mathbf{u}_1 = (0.8507,\ 0.5257)$.
  \item Reconstruction $\hat\vx = \sum_j z_j\mathbf{u}_j$ has mean squared error
        equal to the discarded eigenvalues, so keeping variance and minimising
        reconstruction error are the same goal.
  \item Choose $k$ by variance explained: $21$ of the digits' $64$ components keep
        $90\%$. Two components already separate the digits $0$, $1$ and $7$.
  \item PCA minimises perpendicular, not vertical, distances: it is not regression.
        It ignores labels, so it is a poor cure for overfitting --- regularization
        beat every PCA model on the digits.
  \item t-SNE and UMAP are non-linear maps for looking at data; their inter-cluster
        distances should not be read as measurements.
\end{itemize}

\section*{Exercises}
\addcontentsline{toc}{section}{Exercises}

\exercisesection{Short Questions}

\sq Give three uses of dimensionality reduction.

\sq Why must data be centred before PCA, and when must it be scaled?

\sq What does the eigenvalue of a principal component measure?

\sq Why are the principal components orthogonal?

\sq Define the variance explained by the first $k$ components.

\sq How is an example reconstructed from its first $k$ components?

\sq How does PCA differ from linear regression?

\sq Why is PCA a poor remedy for overfitting?

\sq Why may two implementations of PCA return components of opposite sign?

\sq What do t-SNE and UMAP do that PCA cannot, and what should not be read from
their plots?

\exercisesection{Analytical Problems}

\ap Derive the first principal component by maximising $\mathbf{u}^{\mathsf T}
\Sigma\mathbf{u}$ subject to $\mathbf{u}^{\mathsf T}\mathbf{u} = 1$, and show that
the maximum variance is the largest eigenvalue.

\ap Carry out PCA by hand on the points $(3,1)$, $(1,1)$, $(-1,-1)$, $(-3,-1)$:
compute $\Sigma$, its eigenvalues, the first component, the projections and the
variance explained.

\ap Show that the mean squared reconstruction error from $k$ components is
$\sum_{j>k}\lambda_j$.

\ap Show that if $\Xmat = \mathbf{U}\mathbf{S}\mathbf{V}^{\mathsf T}$ is the SVD of
centred data, the eigenvectors of $\Sigma$ are the columns of $\mathbf{V}$ and the
eigenvalues are $s_j^2/m$.

\ap Find the slope of the least-squares regression of $x_2$ on $x_1$, of $x_1$ on
$x_2$, and of the first principal component, for the five points of
\cref{sec:pca-algorithm}. Explain why the three differ.

\ap A dataset has eigenvalues $4.0$, $2.5$, $0.9$, $0.4$, $0.2$. How many
components keep at least $90\%$ of the variance, and what is the mean squared
reconstruction error with that many?

\exercisesection{Design Problems}

\dpr Survey data has $40$ numeric answers from $2000$ respondents. Design a
PCA-based visualisation: preprocessing, the choice of $k$, and how you would explain
the axes to a non-technical audience.

\dpr A sensor network records $500$ correlated readings per second. Design a PCA
compression scheme, including how you would choose $k$ and how you would detect when
the components need to be recomputed.

\dpr A colleague proposes PCA to $10$ components before training a classifier
``to avoid overfitting''. Design the experiment that tests the proposal against
regularization, and say what result would justify it.

\dpr Design an exploratory analysis of an unfamiliar dataset of $10{,}000$ examples
with $100$ features, using PCA and t-SNE together, and say which conclusions each
can and cannot support.

\exercisesection{Programming Exercises}

\pe Implement PCA with \code{np.linalg.eigh} and with \code{np.linalg.svd}, and
confirm that both reproduce the five-point example.

\pe Reproduce the digits table of \cref{sec:reconstruction}, the reconstructions of
\cref{fig:recon} and the cumulative curve of \cref{fig:cumvar}.

\pe Reproduce \cref{fig:digits-pca}, then do the same for all ten digits, and
compare with a t-SNE map of the same images.

\pe Reproduce the PCA-versus-regularization table, and add a curve of test accuracy
against $k$ from $1$ to $64$.

\exercisesection{Self-Test}

\qz The first principal component is the eigenvector of the covariance matrix
with: (a)~the smallest eigenvalue; (b)~the largest eigenvalue; (c)~eigenvalue $1$;
(d)~any eigenvalue.

\qz The variance of the projections onto a principal component equals:
(a)~$1$; (b)~its eigenvalue; (c)~the number of features; (d)~zero.

\qz Before PCA the data must be: (a)~labelled; (b)~centred; (c)~clustered;
(d)~binary.

\qz The mean squared reconstruction error from $k$ components equals:
(a)~$\lambda_1$; (b)~the sum of the discarded eigenvalues; (c)~$k$;
(d)~zero always.

\qz PCA minimises: (a)~vertical distances; (b)~perpendicular distances;
(c)~horizontal distances; (d)~classification error.

\qz To control overfitting, prefer: (a)~PCA; (b)~regularization;
(c)~more features; (d)~t-SNE.

\qz In a t-SNE plot, the distances between clusters: (a)~are exact; (b)~carry
little meaning; (c)~equal PCA distances; (d)~are always zero.

\furtherreading

\textcite{james2021} §12.2 is the primary reading: it develops principal components
both as directions of maximum variance and as the closest low-dimensional surface,
discusses the proportion of variance explained and the scree plot, and insists on
scaling with a worked example of what goes wrong without it. \textcite{bishop2006}
§12.1 derives PCA by both routes --- maximum variance with a Lagrange multiplier,
exactly as in \cref{sec:pca-derivation}, and minimum reconstruction error --- and
applies it to digit images with reconstructions like those of \cref{fig:recon}.

\bigskip
\noindent\emph{Next:} \Cref{ch:ch17} opens Part~VII, reinforcement learning: an
agent that learns not from labels or structure but from rewards.
